\section{Robustness}
\label{sec:robustness}

\subsection{Threshold and Specification Sensitivity}
\label{sec:rob_threshold}

The headline price gap and its detection-regime gradient survive
reasonable variations in the design choices that define the
screen. The IQR multiplier governs the bright-line cutoff that
separates frequent losers from the broader always-loser stratum.
Under five multipliers from $1.0\!\times$ to $3.0\!\times$
(Table~\ref{tab:threshold}), the price coefficient is positive
and significant in every row, declining monotonically from
$\valThreshMultLow$ at $1.0\!\times$ ($\valThreshMultLowN$ firms)
to $\valThreshMultHigh$ at $3.0\!\times$ ($\valThreshMultHighN$
firms). The decline is a
feature, not a defect: as the cutoff tightens, the binary
indicator filters fewer firms with high loss-intensity, and
\S\ref{sec:mech_horse_race} establishes that the underlying
primitive is the continuous $\log(1+\text{tenders\_count})$
statistic. Standard-error clustering does not change the point
estimate; under item-level clustering (baseline,
SE $\valClusterItem$), procuring-unit-level clustering
(SE $\valClusterPBU$, tighter), and two-way
\citet{cameron2011robust} clustering (SE $\valClusterTwoway$,
looser), the coefficient is significant at $p<0.01$
(Table~\ref{tab:clustering}). The continuous score discriminates
the same cobidder population at AUC $\valAUClogtc$
($95\%$ CI $\valAUClogtcCI$) versus $\valAUCFLBinSameSample$
($\valAUCFLBinSameSampleCI$) for FL14
(DeLong $Z = \valDeLongZ$, $p = \valDeLongP$); standalone price
coefficients of $+0.0653$ ($p<0.001$, FL14 binary) and
$+0.0188$ ($p<0.001$, continuous) point in the same direction
(Table~\ref{tab:horse_race_v14}). The headline price gap is
neither a threshold artifact nor a binary-discretization artifact.

\subsection{Identification Audits}
\label{sec:rob_identification}

Four audits calibrate the descriptive scope of $\beta$ under
specific identification concerns. The
\citet{cinelli2020making} robustness value
$RV_{q=1} = \valRVqOne$ implies an unobserved confounder would
need to explain at least $\valRVqOne$ of the residual variation
in both frequent-loser presence and log price to drive the
coefficient to zero; \citet{oster2019unobservable} bounds
agree directionally with $\hat\delta = \valOsterDelta$ on the
full sample (Table~\ref{tab:oster_delta}). The pregão subsample
yields $\hat\delta = 634.00$; the convite subsample carries a
smaller margin ($\hat\delta = 43.62$), consistent with
\S\ref{sec:mech_modal}'s finding that the construct's
signal-to-noise is highest in pregão environments.

The strict-overlap matching exercise asks a different
identification question. Restricting comparisons to cells in
which frequent-loser items and non-frequent-loser items
genuinely overlap on observables returns
$\widehat\beta^{\text{ov}}=\valMatchOverlapCoef$
($N=\valBetaOverlapN$); propensity-score-trimmed matching returns
$\valMatchPSCoef$. The contrast with the broad-sample
$\widehat\beta = \valBetaBroad$ is decomposed in
\S\ref{sec:results_overlap}: the unweighted overlap regression
returns $\valBetaUnwOverlap$ (only $\valDroppedItemsShare$ of
treated items are strictly dropped), the segment-level
decomposition locates the broad-sample positive in tender-value
Q4 ($\valSRQfourBroad$ robust to all three specifications) and
the negatives in Q1--Q3 ($\valSRQoneBroad$ to
$\valSRQthreeBroad$, also robust), and trim sensitivity
\emph{strengthens} the negative ATT as the most heavily weighted
cells are removed (from $\valSRTrimNone$ to $\valSRTrimTen$ at
the top decile). Both signs describe the same population from
different design angles; we report the four specifications
jointly so the reader can locate the construct's descriptive
claim within the appropriate scope.

The leakage audit (Table~\ref{tab:leakage_audit}) addresses
tautological leakage from cobidders being by construction
elevated-participation always-losers. Holding the score fixed
and swapping the label to direct CADE defendants returns
$\valLeakDirectAUC$, random discrimination that confirms the
construct's structural scope (cover bidders, not ringleaders).
Withholding cobidder firms from score construction within five
cross-validation folds and re-evaluating on items containing
held-out cobidders returns AUC $\valLeakCVAUC$
($95\%$ CI $\valLeakCVAUCCI$); under temporal holdout
(train $2009$--$2016$, test $2017$--$2019$) AUC is
$\valAUCFLfirmTemp$ ($95\%$ CI $\valAUCFLfirmTempCI$). The
$0.10$--$0.13$ drop from the in-sample reference is the
pure-leakage component, and we adopt the temporal-holdout AUC of
$\valAUCFLfirmTemp$ as the conservative discriminating reference.
A placebo specification using the supply of sub-threshold
always-losers as instrument confirms the construct's price
imprint draws on the above-threshold cover-bidder population, not
on generic always-loser supply (Online Appendix~C,
Table~\ref{tab:iv_placebo}).

\subsection{Temporal Holdout: Generalization of the Screening Statistic}
\label{sec:rob_operational}

In-sample evaluation confounds classification with prediction:
items in the late panel contribute both to score construction and
to the labels the score is evaluated against. Constructing the
score on $2009$--$2016$ participation only and evaluating on
items in $2017$--$2019$ yields a firm-level AUC of
$\valAUCFLfirmTemp$ ($95\%$ CI $\valAUCFLfirmTempCI$) against the
$\valCobidders$ adjudicated CADE-cobidder labels---the
generalization reference for the screening statistic throughout
the paper. Operational metrics
(Table~\ref{tab:operational_metrics}) show in-sample precision
overstates by approximately $50\%$ at top-$500$
($\valOpInsamplePrecFiveHund$ in-sample versus
$\valOpPrecFiveHund$ holdout) because roughly
$\valOpInflationShare$ of the top-$500$ in-sample ranking comes
from $2017$--$2019$ participation, after CADE adjudications were
already underway for some cartels. We disclose the gap explicitly
and report headline metrics on the holdout column. At top-$1{,}000$
the screen flags adjudicated cobidder firms with recall
$\valOpRecallThou$ under temporal holdout, and the architectural
extension of the temporal-holdout audit to the
sequential-gatekeeper rule appears in
Table~\ref{tab:architecture_gatekeeper_th} (\S\ref{sec:forensic}).
